\section{Introduction}
\label{sec:introduction}

Fast-charging protocol design must balance short charging time against
temperature rise and degradation-related behavior. Large currents can reduce
the time required to reach a target state of charge (SOC), but the resulting
electrochemical and thermal response depends nonlinearly on the cell state and
protocol sequence~\cite{ref2,ref3,ref23}. We therefore study three-stage
constant-current protocols under encoded voltage, temperature, SOC, and
decision-variable limits, with charging time, peak temperature rise, and an
uncalibrated protocol-level degradation proxy as three minimized objectives.

Each candidate requires a coupled simulation or physical experiment, making
BO attractive for its probabilistic surrogate and acquisition-driven
evaluation policy~\cite{ref4}. Prior battery studies use constrained and
multiobjective BO for charging-protocol search~\cite{ref7,ref8,ref15,ref18,
ref19}, but early decisions still have little task-specific data. LLM-assisted
methods can supply initialization or search preferences~\cite{ref9,ref10,
ref17}; because these outputs are neither measurements nor calibrated
posterior evidence~\cite{ref16}, they must remain separate from simulator
observations.

LLMBO-MO enforces this separation with a screened warm-start portfolio and a
weight-conditioned point or region preference. Depending on configuration,
the preference augments a plain expected-improvement (EI) search or modifies
the acquisition mean through posterior covariance; only simulator evaluations
create objective vectors. Figure~\ref{fig:framework} summarizes the workflow.

Five paired seeds and 56 calls provide a fixed-budget comparison on two cell
parameterizations; all differences are reported descriptively.

The contributions are threefold:
\begin{enumerate}
\item We formulate a simulation-based, three-objective battery fast-charging
problem and make its implementation boundary explicit, including the
uncalibrated status of the degradation proxy.
\item We present an LLM-augmented ParEGO workflow with a deterministic
warm-start portfolio, weight-conditioned preferences, and an
implementation-faithful posterior-covariance mean-lift construction whose
covariance is unchanged and whose mean shift is budgeted.
\item We provide a paired, fixed-budget evaluation on two PyBaMM
parameterizations together with matched trajectories, protocol inspection,
runtime, and mechanism telemetry.
\end{enumerate}
